\section{Background and Problem Setup}
\label{sec:problem}

\subsection{Dynamic borrowing setup}

Let
\[
D_c = \{(x_{ci}, y_{ci})\}_{i=1}^{n_c}
\quad \text{and} \quad
D_e = \{(x_{ej}, y_{ej})\}_{j=1}^{n_e}
\]
denote the concurrent and external datasets, respectively.
The inferential goal is always defined with respect to the concurrent
population.
External data can be useful when it is sufficiently compatible with the
concurrent cohort, but should be downweighted or ignored when the two sources
are not exchangeable.

In the stylized simulation study, the target estimand is a scalar
current-study treatment effect $\theta$.
The NPE method receives fixed-length summary statistics computed from
$(D_c, D_e)$ and returns an approximate posterior
\[
q_{\phi}(\theta \mid s(D_c, D_e))
\approx
p(\theta \mid D_c, D_e),
\]
where $s(D_c, D_e)$ contains source-level summaries and mismatch diagnostics.
This preserves the core dynamic borrowing goal while replacing repeated
per-dataset Bayesian fitting with amortized posterior inference.

\subsection{Exchangeability regimes}

The simulation study is organized around six regimes designed to stress-test
borrowing methods under distinct forms of source mismatch:

\begin{center}
\begin{tabular}{>{\raggedright\arraybackslash}p{0.14\linewidth}
                p{0.78\linewidth}}
\toprule
Scenario & Interpretation \\
\midrule
SC1 & \textbf{All exchangeable.}
      The external and concurrent cohorts have similar covariate distributions
      and the same outcome mechanism. \\[2pt]
SC2 & \textbf{Partial covariate shift.}
      Moderate differences in covariate distributions with a similar outcome
      mechanism. \\[2pt]
SC3 & \textbf{Partial outcome drift.}
      Covariate distributions are similar, but the outcome mechanism differs
      moderately. \\[2pt]
SC4 & \textbf{Strong covariate shift.}
      The external cohort differs substantially in case mix. \\[2pt]
SC5 & \textbf{Strong outcome drift.}
      Covariate distributions are similar, but the outcome mechanism differs
      materially. \\[2pt]
SC6 & \textbf{Joint non-exchangeability.}
      Both the covariate distribution and the outcome mechanism differ. \\
\bottomrule
\end{tabular}
\end{center}

These regimes deliberately separate settings in which mismatch is mainly
covariate-driven from settings in which it is mainly outcome-driven.
They provide a common evaluation grid for both classical borrowing methods and
the amortized NPE estimator.

\subsection{Summary-based neural posterior estimation}

The NPE model operates on a fixed-length summary vector
\[
s(D_c, D_e)
=
\bigl[
\mu_c,\sigma_c,\mu_e,\sigma_e,\Delta_{\mu},\Delta_{\beta},\text{size terms},
\text{overlap terms}
\bigr],
\]
where the exact contents are determined by the method implementation.
In the current baseline NPE model, the summary includes source-level means and
standard deviations, outcome summaries, simple regression-based coefficients,
and overlap-inspired diagnostics derived from source-membership propensity
scores.

Given simulated pairs $(\theta, D_c, D_e)$, a neural network is trained to
approximate the posterior distribution over $\theta$.
The advantage of this setup is computational: training is done once offline,
and posterior inference for a new current/external pair requires only a single
forward pass through the network.
This makes the approach attractive in settings where many candidate external
sources might be screened repeatedly.

